He could not lace a boot.

That was the thing that got him, on the fourteenth morning, sitting on the edge of a cot in a room off the armoury with one boot on and the other in his lap. Not the sword. He had known about the sword since the third day, when Sam had told him straight out that he would not know from one morning to the next whether the hand would hold one, and he had taken that the way a man takes a thing he has already half worked out for himself.

The boot he had not thought about.

You needed two hands for a lace and you needed the fingers on both of them, and the two on his right that Sam said would never do much did not do much, and the other three worked in a general sort of way and would not do the small work at all. He got the lace through the eyelet. He could not make the loop. He tried it eleven or twelve times and got angry at about the sixth and went on anyway, because going on anyway was the only method he had ever had for anything, and at the end of it his hand was shaking hard enough that he had to put it under his arm and wait.

Then he went out with the boot loose and told himself he would find somebody later.

He did not find anybody. He wore it loose for two days.

\sceneBreak

The trouble was that there was nothing for him to do, and it had taken him until about the tenth day to work out that this was a trouble and not a mercy.

Everyone had a thing. Tyrion had the city, in the way a man has a horse that is running away with him. Sansa had the North and whatever the North was now. Sam had four hundred people on a floor. Davos had the food, Ser Roland had the ground, the Greyjoy woman had the ships, and Grey Worm had eight thousand men who wanted telling.

Jon had killed the Night King.

That was not a thing you could get up and do in the morning. It had taken about four seconds and it was finished and there was no second one to be done, and he had discovered over a fortnight that being the man who had done it was not an occupation but a condition, like weather.

So on the fourteenth day he went down to the water lines and asked for work.

\sceneBreak

It went wrong within a quarter of an hour and he did not understand why for some time.

The line at the Cobbler’s Square well ran three hundred people long and there were four men on the winch and a fifth turning people away when the bucket came up muddy, and Jon put himself on the winch and got about eleven turns into it before the man beside him stopped.

Then the man on the other side stopped.

Then the line stopped, which he had not expected at all, because a line of three hundred people waiting for water does not stop for anything, and Jon straightened up with the winch handle in his left hand and found that a considerable number of people were looking at him and that the looking had a quality to it he had no name for.

“You’ll want the other side, my lord,” said the man beside him. “Left-hand crank. It’s easier on you.”

“I’m all right on this one.”

“Aye,” the man said. “But you’ll want the other side.”

And he moved. He picked up his coat and went round the winch and stood on the far side of it and left Jon the left-hand crank, and did it politely and immediately and would not be argued out of it, and the whole line watched him do it and nobody said a word, and the line started up again.

Jon turned the crank for four hours. Nobody would let him carry anything.

\sceneBreak

By the third hour they had started bringing things.

That was how he came to understand it, in the end — not from the looking, which he had put down to the hand, but from the objects. A woman came out of the line and put a rope of dried figs on the wall beside him and went back to her place without saying anything. Somebody left a knife, a good one, a fisherman’s, with the handle wrapped. A boy of about ten came right up and gave him a bird’s skull and stood there with his hands behind his back until Jon took it, and then went off at a run.

He tried giving the figs back and the woman would not have them.

At the fourth hour a man came up the line out of turn with a child on his hip, a girl of perhaps four with her eyes gummed shut and her breath going wrong, and he did not ask for water.

“Just the hand,” he said. “That’s all. Just put it on her.”

“I can’t do anything for her.”

“I’m not asking you to do anything.” The man’s voice was quite level and quite reasonable and he had clearly been rehearsing it. “I’ve had her to the maester. There’s forty in front of her and he’s honest with me and I’ve no complaint of him. I’m asking you to put your hand on her head, and then I’ll go, and I’ll not trouble you again.”

There were perhaps two hundred people watching.

Jon stood there with a winch handle in his left hand and worked through it, slowly, the way he worked through everything, and he could not find the end of it that let him say no. If he said no he was refusing a man with a sick child in front of two hundred people, and the man would go away and tell it, and there was no version of that telling in which Jon came out as a man who had been honest rather than a man who had withheld something.

So he put his hand on her head.

It was the ruined one, because the man had asked for the hand and there was only one hand he could have meant, and he put it on the girl’s head for about three seconds and took it off, and the man said thank you and went, and Jon felt nothing whatsoever except that the child was too hot, which he already knew from looking at her.

Two hundred people watched him do it.

He understood by the fifth hour that he had made an enormous mistake and he could not for the life of him identify what else he could have done.

\sceneBreak

Sansa found him at the seventh hour and told him about the North.

“Nothing,” she said. “Not from White Harbor. Not from anywhere. Four birds out and none back, and Tyrion’s the same, and Lord Davos has had one ship north of Duskendale in a fortnight and it turned at Maidenpool because of the ice.”

“They’ll be shut in. Winter and the roads and no ravens — half those keeps have one maester and he’ll be dead or he’ll be walking.”

“Yes.”

“That’s not the same as nobody being there.”

“No,” Sansa said. “It isn’t. I’ve told myself that four times a day for eleven days and I’ve stopped finding it convincing, and I’d rather you knew that I’ve stopped.”

She had a way of putting a thing down in front of him and waiting, and he had never once in his life been able to do anything with it except pick it up.

“Then somebody has to go and look.”

“Somebody will. Not before the thaw and not on a ship with grain in it, and it wants a man who can be spared.” She was looking at him carefully. “I’d rather it wasn’t you.”

“Why?”

Sansa did not answer that directly, which he noticed and did not pursue.

“There were six hundred and forty Northmen in this city this morning,” she said instead. “There are six hundred and thirty-two tonight. Two in the arch, four of the flux, one had his skull opened by a Reach man over a woman and one walked out of the Iron Gate on the eleventh day and hasn’t come back.” She smoothed her skirt. “That last one is the number I watch. If it goes to twenty, I’ve got a problem I can’t fix by standing in front of them.”

“They’ll hold. They’re Northmen and they’ll hold.”

“They’re Northmen who won a war eight hundred miles from home and have had no word from it in a fortnight,” Sansa said, “and they are being fed worse than the Unsullied, and every single one of them can count.” She stood. “You’re wanted at the pits tomorrow, by the way. Ser Roland asked whether you’d stand at the reading of the names.”

“For what purpose?”

“For the purpose of you standing there.”

“That’s not a reason.”

“Jon.” She stopped at the door. “It is the only reason anybody has for anything at the moment, and I’d take it and be grateful. There are people going into that pit whose families think it’s a hole. If you stand at the head of it for an hour they’ll think it’s something else, and it will cost you an hour.”

He said he would think about it, which meant he would do it, and they both knew that.

\sceneBreak

He had not been back to see Daenerys.

He had gone once, on the second night. She had been asleep on the floor of the Great Hall in a row with a hundred others and he had sat down on the stone beside her for about an hour and had not woken her, and then he had got up and gone, and he had not been back in twelve days.

He had four reasons and they were all true.

The hand — Sam had him back twice a day for the first week and there had been three days when he could not have crossed the yard. Her state — they had told him she could not sit up, and a man standing over a woman who cannot sit up is doing something to her whether he means to or not. The thing about being watched — there were four hundred people on that floor and every single one of them would have seen it, and there were already enough stories.

The fourth one he did not examine and could have stated in one sentence if he had ever tried, which was that on the flat, before she went up, she had said \textit{when I come down, there’s something I want to tell you}, and she had come down, and neither of them had said anything, and every day that passed made the not-saying of it into a larger and more particular object that would have to be got round before anything else could be said at all.

He thought about it about six times a day and each time arrived at the conclusion that he would go tomorrow, and each time the conclusion felt like a decision, which was the part that let it go on working.

He heard on the fourteenth day, secondhand and four days late, that she had been moved out of the Great Hall. Nobody had told him — there was no one whose task that was, and no list of people owed word of anything, which was less an oversight than a description of the whole government of Westeros just now. He was annoyed about it in the flat, worn-down way he was annoyed about most things that fortnight, and did not chase it further.

\sceneBreak

The Golden Company man was on the tenth of the water carts and Jon nearly missed him entirely.

There were about nine hundred of them left and they were being fed and worked, which was as much as anybody could manage, and this one was fifty-odd and had a Westerosi accent under the sellsword in it and had been shifting barrels since before dawn. He straightened up when Jon came past and looked at him for slightly too long, and Jon, who had been looked at by a great many people that day, stopped.

“Say it, then.”

“Beg pardon, my lord.” The man was not embarrassed. “It’s only I’d heard about the hand and I’d not seen it.”

“Then look at it.” Jon held it up. It was still wrapped to the wrist and the two fingers sat wrong even under the linen. “It’s not much of a thing to see. It went in a fire and it came out worse and there’s not anything in it beyond that.”

“Aye,” the sellsword said. “Fire took you, then.”

“It did.”

“Right.” The man nodded slowly, twice, the way a man does when a question he has been carrying for some while has finally been settled, and he did not appear disappointed and he did not appear pleased. “That’s answered, then. Thank you, my lord.”

“That’s answered? What’s answered?”

“Nothing, my lord.” He bent to the barrels. “Only there’s been talk and there’s always talk. Nobody’s business.”

Jon stood there a moment longer.

He knew the man was not telling him something and he could see it plainly, and he also knew that there were nine hundred of these men in the city with nothing to their name and every reason to be careful with a lord, and that a sellsword being cagey with the man who had taken his employer’s city was not a mystery, it was a Tuesday.

So he let it go.

He would remember the exchange perfectly, word for word, some months afterward, in a room with three other people in it, and would understand then exactly what had been asked of him and exactly what answer he had given.

\sceneBreak

He worked out what to do about it that evening and was pleased with himself, which ought to have warned him.

The trouble, as he saw it, was that people had got hold of a version of the thing on the flat that was not the true one. He had heard four of them by the fourteenth day. In one of them he had walked into the fire. In one he had held the blade for a count of a hundred. In one — this was the one he had heard from a Northman, of all people, a man from Barrowton who had been thirty feet away and ought to have known better — he had picked the dagger out of the wildfire and it had not burned him at all until afterward, when the thing was dead, and only then had the hand gone.

That last one had a shape to it he did not like.

And the answer, plainly, was the truth, because the truth was the least impressive account available and he happened to have it. He had watched steel do nothing for two days. He had seen a blade in a fire and put his hand in before he had finished the thought. There had been no deciding, no walking, no count of a hundred; there had been a man doing the only thing left in front of him, badly, with no time, and getting lucky, and screaming a great deal afterward. Any man in that line would have done it. Most of them would have done it better.

He would simply say so. Out loud, to as many people as would hear it, until it stopped being a story worth telling.

\sceneBreak

He said it the next morning at the pits, before the reading of the names, to somewhere between two and three thousand people.

He was not good at it and he knew he was not good at it, and he thought at the time that this was probably an advantage. He said that the blade was in the fire because a girl had dropped it. He said he had not decided anything, that there had not been time to decide anything, and that a man who has not decided a thing does not deserve to be praised for it. He said that Ser Brienne of Tarth had gone into an archway eight days ago having decided it completely, in full knowledge, and was in the pit behind him, and that if anyone wanted a thing to tell their children about then that was the one and he would be obliged if they would use it instead.

He said it plainly and he did not dress it and it took about two minutes.

Then he stood at the head of the pit for an hour while the names were read, the way Sansa had asked, and it was very cold and very slow and there were four hundred and eleven names and the reading of them was the worst hour he had spent since the flat.

Afterward a great many people came past and a number of them touched his sleeve, and he let them, because there was no way of not letting them that did not involve pushing somebody, and he went back up the hill in the dark feeling better than he had felt in a fortnight.

\textit{That’s the end of that,} he thought. \textit{Say it plain and say it once and it stops.}

\sceneBreak

Two streets down, a woman who had been standing near the front went home to Rope Lane and told her sister what the King in the North had said at the pits, and her sister told it in the market the next morning.

It was a good telling. It had the girl and the dropped knife in it, and it had the burning, and it had him standing up in front of the whole city and saying he had done nothing worth the telling and to talk about the dead knight instead.

By the third time it was told that morning, that last part had gone to the front, where it worked better.

By the end of the week nobody on the Hook could have told you who Brienne of Tarth was, and everybody could have told you that the man who killed the Long Night had stood at the mouth of the pit with his hand in rags and said he was nobody.